%!TEX TS-program = xelatex
%!TEX encoding = UTF-8 Unicode
\documentclass[manuscript,screen,review]{acmart}

\setcopyright{none}
\copyrightyear{2026}
\acmYear{2026}
\acmJournal{TOSEM}
\settopmatter{printacmref=true, printccs=true, printfolios=true}
\setcitestyle{numbers,sort&compress,square}

\usepackage{amssymb,amsthm}
\usepackage{mathtools}
\usepackage{calc}
\usepackage{enumitem}
\usepackage{longtable}
\usepackage{array}
\usepackage{listings}
\usepackage{pifont}

\setlength{\emergencystretch}{3em}
\lstset{
  breaklines=true,
  breakatwhitespace=true,
  basicstyle=\ttfamily\small,
  frame=none,
  columns=fullflexible
}

\title[Supplementary Material]{Supplementary Material for ``A Semantic Mutation Metric for Metamorphic-Relation Adequacy in Scientific Computing Programs''}

\author{Meng Li}
\orcid{0000-0002-1074-1502}
\email{mlemon@usc.edu.cn}
\affiliation{%
  \institution{School of Computing, University of South China}
  \city{Hengyang}
  \country{China}
}
\affiliation{%
  \institution{Hunan Engineering Research Center of Software Evaluation and Testing for Intellectual Equipment}
  \city{Hengyang}
  \country{China}
}
\affiliation{%
  \institution{CNNC Key Laboratory on High Trusted Computing}
  \city{Hengyang}
  \country{China}
}

\author{Xiaohua Yang}
\orcid{0000-0002-2977-1787}
\email{xiaohua1963@foxmail.com}
\affiliation{%
  \institution{School of Computing, University of South China}
  \city{Hengyang}
  \country{China}
}

\author{Jie Liu}
\orcid{0009-0008-1970-8347}
\email{jieliu@usc.edu.cn}
\affiliation{%
  \institution{School of Computing, University of South China}
  \city{Hengyang}
  \country{China}
}

\author{Shiyu Yan}
\orcid{0000-0001-7626-5185}
\email{yanshiyu@usc.edu.cn}
\affiliation{%
  \institution{School of Computing, University of South China}
  \city{Hengyang}
  \country{China}
}

\begin{CCSXML}
<ccs2012>
 <concept>
  <concept_id>10011007.10011074.10011099.10011102</concept_id>
  <concept_desc>Software and its engineering~Software testing and debugging</concept_desc>
  <concept_significance>500</concept_significance>
 </concept>
 <concept>
  <concept_id>10011007.10011074.10011099</concept_id>
  <concept_desc>Software and its engineering~Software verification and validation</concept_desc>
  <concept_significance>300</concept_significance>
 </concept>
</ccs2012>
\end{CCSXML}

\ccsdesc[500]{Software and its engineering~Software testing and debugging}
\ccsdesc[300]{Software and its engineering~Software verification and validation}

\keywords{metamorphic testing, mutation testing, semantic mutation testing, metamorphic relation adequacy, metamorphic relation assessment, scientific computing}

\begin{document}

\maketitle

\appendix

\section{Notation and Operator Catalogue}\label{a.-notation-and-operator-catalogue}

\subsection{Complete notation table}\label{a.1-complete-notation-table}

\paragraph{Experimental objects.}
\begin{itemize}
\item PUT set $I = \{\text{A1,A2,A3, B1,B2,B3, C1,C2,C3, D1,D2,D3}\}$, $|I| = 12$.
\item PUT: $S_i$ ($i \in I$).
\item Class mapping: $\mathrm{cls} : I \to \{A, B, C, D\}$ (open, extensible).
\item Valid input distribution: $D_S$, sampling $X_{K_{\mathrm{eq}}} \sim D_S$.
\end{itemize}

\paragraph{Metamorphic testing side.}
\begin{itemize}
\item Metamorphic-relation label set: $\mathcal{R} = \{\text{Cons.}, \text{Mono.}, \text{Conv.}, \text{Traj.}, \text{P-ord.}\}$, $k \in \{1,\ldots,5\}$ (open, extensible).
\item Metamorphic Relation: $\mathrm{MR}_{i,k}$ (from the shared infrastructure); $\mathit{mr} = (r, R) \in \mathrm{MR}_{i,k}$.
\item Automated Verification Pipeline (AVP): $\mathrm{AVP} : \text{Programs} \times \mathrm{MR}_{\text{universe}} \times \mathbb{R}^+ \to \{\text{pass}, \text{fail}\}$; $\mathrm{AVP}(s, \mathit{mr}, \varepsilon_{\mathrm{AVP}}^k)$ invokes the verification method per relation index $k$.
\end{itemize}

\paragraph{Mutation testing side.}
\begin{itemize}
\item Mutation operator family: $\mathrm{MUT} = \{\mathrm{mut}_1, \ldots, \mathrm{mut}_5\}$, $j \in \{1,\ldots,5\}$ (open, extensible); $\mathrm{mut}_1 = \mathrm{mut}_C$, $\mathrm{mut}_2 = \mathrm{mut}_M$, $\mathrm{mut}_3 = \mathrm{mut}_G$, $\mathrm{mut}_4 = \mathrm{mut}_T$, $\mathrm{mut}_5 = \mathrm{mut}_F$.
\item Signature: $\mathrm{mut}_j : \text{Programs} \to 2^{\text{Programs}}$.
\item Mutant: $s' \in \mathrm{mut}_j(S_i)$.
\item Alignment: $\mathrm{align}(j) = j$.
\end{itemize}

\paragraph{Three-state decomposition (fixed).}
\[
\begin{aligned}
\mathrm{mut}_j(S_i)
  &= \mathrm{equiv}_{i,k,j} \cup \mathrm{killed}_{i,k,j}
     \cup \mathrm{survive}_{i,k,j},\\
&\hspace{1.5em}\text{with the three sets disjoint, mutually exclusive, and exhaustive.}
\end{aligned}
\]
\begin{itemize}
\item \emph{equiv determination} (dual conditions):
  (E1) AVP-coherent: $\forall \mathit{mr} \in \mathrm{MR}_{i,k}: \mathrm{AVP}(S_i, \mathit{mr}) = \mathrm{AVP}(s', \mathit{mr})$;
  (E2) Output-equiv: $\forall x \in X_{K_{\mathrm{eq}}} \sim D_S: \|S_i(x) - s'(x)\| \leq \varepsilon_{\mathrm{eq}}$.
\item \emph{killed determination}: $\mathrm{killed}(s', \mathrm{MR}_{i,k}) \Leftrightarrow \exists \mathit{mr} \in \mathrm{MR}_{i,k}: \mathrm{AVP}(S_i, \mathit{mr}) = \text{pass} \wedge \mathrm{AVP}(s', \mathit{mr}) = \text{fail}$.
\end{itemize}

\paragraph{Metric (fixed classical structure).}
\begin{align*}
\mathrm{SMS}_{i,k,j} &:= \frac{|\mathrm{killed}_{i,k,j}|}{|\mathrm{mut}_j(S_i)| - |\mathrm{equiv}_{i,k,j}|} \\
                    &\phantom{:}= \frac{|\mathrm{killed}_{i,k,j}|}{|\mathrm{killed}_{i,k,j}| + |\mathrm{survive}_{i,k,j}|} \in [0, 1].
\end{align*}

\paragraph{LRCA engineering attribution layer (descriptive, not in SMS).}
\begin{itemize}
\item Likely root-cause inventory: $C = \{C_1, \ldots, C_5\}$ (open, extensible). C1 True semantic failure; C2 Tolerance perturbation; C3 OOD; C4 Statistical-assumption violation; C5 Mutator artefact.
\item LRCA output: each $s' \in \mathrm{killed}$ annotated with $\texttt{root\_cause}(s') \in C$.
\item Descriptive quantities: $\texttt{C1\_share}_{i,k,j}$, $\texttt{suspect\_share}_{i,k,j}$.
\end{itemize}

\paragraph{Slicing and cross-class.}
\begin{itemize}
\item Slicing: aligned $j = k$; cross $j \neq k$.
\item Cross-class: $\Delta\mathrm{SMS}_c$, $c \in \{A, B, C, D\}$; $\mathrm{CV}(\Delta\mathrm{SMS}) := \mathrm{std}(\Delta\mathrm{SMS}_c) / |\mathrm{mean}(\Delta\mathrm{SMS}_c)|$.
\end{itemize}

\paragraph{Tolerance and sampling.}
\begin{itemize}
\item Tolerance: $\varepsilon_{\mathrm{AVP}}^k$ (relation-dependent), $\varepsilon_{\mathrm{eq}}$ (equivalence output tolerance).
\item Sampling: $K_{\mathrm{eq}} = 1000$ (E2 sample size), $N = 20$ (statistical replicates).
\end{itemize}

The notation skeleton (11 core concepts + three-state decomposition +
SMS formula + AVP interface) is \textbf{fixed}; the \emph{content} of
MUT, $\mathcal{R}$, C, and cls is open to extension.

\subsection{Root-Cause Decision Tree and Diagnostic Protocol}\label{a.2-lrca-decision-tree-and-three-layer-diagnostic-protocol}

\textbf{Five likely root causes} (open):

{%
\begin{longtable}[]{@{}ll@{}}
\caption{Root-cause decision tree and three-layer diagnostic protocol.}\label{tab:p2-15}\\
\toprule\noalign{}
Code & Meaning \\
\midrule\noalign{}
\endfirsthead
\endhead
\bottomrule\noalign{}
\endlastfoot
C1 & True semantic failure (what SMS seeks) \\
C2 & Numerical tolerance perturbation \\
C3 & Out-of-distribution (OOD) \\
C4 & Statistical-assumption violation \\
C5 & Mutator artefact \\
\end{longtable}
}

\textbf{Three-layer diagnosis with L0 artefact pre-scan:}

\begin{itemize}
\item \textbf{L0 artefact pre-scan.} Double-blind review following the main manuscript's mutant-pool prescreen and equivalence-detection protocol: dual-LLM cross-source review plus 20\% manual sampling.
\item \textbf{L1 tolerance robustness.} $N = 20$ replicates; fail ratio $\geq 0.80$ considered stable.
\item \textbf{L2 OOD triage.} For C/D classes, distinguish valid $D_S^{\text{valid}}$ from OOD.
\item \textbf{L3 statistical baseline.} For Wilcoxon/DTW, IID/stationarity pre-check on PUT samples.
\end{itemize}

\textbf{Decision tree:}

For each $s' \in \mathrm{killed}_{i,k,j}$:
\begin{itemize}
\item \textbf{L1.} fail ratio $< 0.80 \Rightarrow$ C2; otherwise $\Rightarrow$ L2.
\item \textbf{L2} (C/D classes). fail only in OOD $\Rightarrow$ C3; otherwise $\Rightarrow$ L3.
\item \textbf{L3} (B/D classes + Wilcoxon/DTW). assumption violation $\Rightarrow$ C4; otherwise $\Rightarrow$ artefact recheck.
\item \textbf{Artefact recheck.} artefact evidence $\Rightarrow$ C5; otherwise $\Rightarrow$ C1.
\end{itemize}

\textbf{Multi-label priority:} C5 \textgreater{} C4 \textgreater{} C3
\textgreater{} C2 \textgreater{} C1.

\textbf{LRCA threshold calibration (9-grid scan;
\texttt{lrca\_calibration.json}).} Best ood\_band=0.02 with any
tolerance\_multiplier (3.0 / 10.0 / 30.0) gives mean\_C1\_share 0.200 /
H4 12/60 (20.0\%); ood\_band=0.05 (default) or 0.10 gives 0.164 / 10/60
(16.7\%). \texttt{tolerance\_multiplier} has zero impact (L1 tolerance
rarely triggered); \texttt{ood\_band} is the only discriminator.
Calibration ceiling 12/60 = 20\% remains far below 80\%, H4
unattainable on this dataset, not a calibration issue (inherent SNR of
LLM-mutant pools).

\subsection{Verification Protocol and Equivalence-Judgement Trade-Off}\label{a.3-avp-protocol-specification-and-equivalence-judgement-trade-off}

\[
\mathrm{AVP}: \mathrm{Programs} \times \mathrm{MR}_{\mathrm{universe}} \times \mathbb{R}^{+} \to \{\mathrm{pass}, \mathrm{fail}\}.
\]
Per-relation verification: $\mathrm{Cons.}$ tolerance equality $|\mathrm{LHS} -
\mathrm{RHS}| \leq \varepsilon$; $\mathrm{Mono.}$ / $\mathrm{P-ord.}$ Wilcoxon signed-rank $\alpha=0.05$; $\mathrm{Conv.}$
convergence-order + asymptotic residual ratio; Traj. DTW distance
threshold $\varepsilon_{\mathrm{DTW}}$. In NOETHER's taxonomy
\citep{noether2026} these are $m_{\mathrm{inv}}$, $m_{\mathrm{mono}}$,
$m_{\mathrm{cmp}}$, $m_{\mathrm{conv}}$, and $m_{\mathrm{dyn}}$; note that a
P-ord ($m_{\mathrm{cmp}}$) instance reducing to a pure asymptotic
residual-ratio rate check shares its acceptance machinery with the Conv
($m_{\mathrm{conv}}$, limit block $\mathcal{L}^{*}$) checker, so such
rate-only instances are $\mathcal{L}^{*}$-adjacent while the registered
stratum stays partial-order. AVP version = upstream infrastructure commit
hash; our package embeds AVP source.

\textbf{Three-candidate equivalence trade-off.}

{%
\begin{longtable}[]{@{}
  >{\raggedright\arraybackslash}p{(\linewidth - 6\tabcolsep) * \real{0.2500}}
  >{\raggedright\arraybackslash}p{(\linewidth - 6\tabcolsep) * \real{0.2500}}
  >{\raggedright\arraybackslash}p{(\linewidth - 6\tabcolsep) * \real{0.2500}}
  >{\raggedright\arraybackslash}p{(\linewidth - 6\tabcolsep) * \real{0.2500}}@{}}
\caption{Verification protocol and equivalence-judgement trade-off.}\label{tab:p2-16}\\
\toprule\noalign{}
\begin{minipage}[b]{\linewidth}\raggedright
Determination
\end{minipage} & \begin{minipage}[b]{\linewidth}\raggedright
False-positive
\end{minipage} & \begin{minipage}[b]{\linewidth}\raggedright
False-negative
\end{minipage} & \begin{minipage}[b]{\linewidth}\raggedright
SMS bias
\end{minipage} \\
\midrule\noalign{}
\endfirsthead
\endhead
\bottomrule\noalign{}
\endlastfoot
E1 alone & numerical coincidence + AVP-coherent & E1 false $\to$ non-equiv &
Low (smaller denominator) \\
E2 alone & output match, MR differs (rare) & K\_eq miss, full-space
equal & High \\
\textbf{E1 $\wedge$ E2} & both mis-pass simultaneously & E1 $\vee$ E2 false $\to$ non-equiv
& Slightly high (conservative) \\
\end{longtable}
}

E1 alone is fooled by insufficient AVP coverage; E2 alone by numerical
coincidence on $K_{\mathrm{eq}}$. E1 $\wedge$ E2 is the conservative complete
instantiation; under $L_{\mathrm{equiv}}$ it reduces almost-everywhere to
classical bitwise equivalence (Lemma G.1). Counter-examples: E2 passes
but E1 does not (rare; mutant coincides at K\_eq but deviates at MR
trigger points $\to$ non-equivalent); E1 passes but E2 does not (common;
consistent within MR but numerical drift \textgreater{} $\varepsilon_{\mathrm{eq}} \to$
non-equivalent).

% [thematic break removed]

\section{Experimental Subjects and Mutation-Operator Specialisations}\label{b.-experimental-subjects-and-mutation-operator-specialisations}

\subsection{Program Selection Rationale}\label{b.1-put-selection-rationale-detailed-coverage-argument}

\textbf{(a) Library-stack coverage.} numpy 2.4.4 (A1-B3 linear algebra /
arrays), scipy 1.17.1 (A1 integrate, A2 linalg, B2 stats), scikit-learn
1.8.0 (C1-D3 surrogate / ML), Python scientific-computing's de facto
foundation stack (PyPI top-50, 2026-04). Uncovered: GPU / distributed
(JAX, CuPy, dask), domain-specific (BioPython, Astropy, RDKit). Left to
future work under the representativeness threat.

\textbf{(b) Mathematical-structure coverage.} ODE/PDE (A1, A3), direct
linear algebra (A2), Bayesian analytic + MCMC + Monte Carlo (B1-B3),
kernel + orthogonal polynomial + NN surrogate (C1-C3), convex
optimisation + backprop + max-likelihood (D1-D3), covers 8 of 12
chapters in Numerical Recipes \citep{press2007numerical}. Uncovered: advanced
PDE solvers (FEM, FV, spectral); FFT; interior-point/trust-region
optimisation; symbolic/CAS. Left to future work.

\textbf{(c) Comparison with benchmarks.} DeepCrime (Humbatova et
al.~2021): semantic class-D topical overlap on ML kernels. Defects4J (Just et
al.~2014): no overlap, only mutation-as-fault-proxy reference. mutmut /
cosmic-ray demos: general Python; the semantic operators extend the scientific-computing focus.

\textbf{(d) Scale vs representativeness.} Each PUT 50-400 LOC; signature
standardised to \texttt{program(x:\ float)\ $\to$\ float}. Substantively
smaller than industrial code (typically 1-10 KLOC).

\textbf{Limitation.} The \texttt{float\ $\to$\ float} signature is a
substantive constraint, not engineering trade-off. Industrial inputs are
typically high-dimensional (CFD grids, MD particles, FEM matrices);
scalarised PUTs upper-bound mutant semantic complexity and may
systematically under-estimate SMS and cross-class differences on
industrial PUTs (declared in the main manuscript's threats-to-validity section;
future work validates on industrial-grade PUTs).

\subsection{Systematic vs incidental}\label{b.1.5-systematic-vs-incidental}

Syntactic tools may \emph{occasionally} hit the semantic-mutation criteria in the main manuscript (a)
(cross-function-boundary) or (c) (algorithm class), but occasionality
undermines two engineering functions of semantic mutation. \textbf{(a)
Deepening source-code understanding.} Designing OS
\texttt{np.linalg.det(M)\ $\to$\ np.sum(np.diag(M))} requires knowing the
two are equivalent on diagonal matrices but not on general matrices
(\texttt{det\ =\ $\prod$\ eigvals} vs
\texttt{sum(diag)\ =\ trace\ =\ $\sum$\ eigvals}); syntactic tool AST
traversal does not require this understanding even when similar
replacements appear. \textbf{(b) Revealing deep faults.}
Domain-semantic errors (physical-constant, unit conversion, boundary
conditions, hyperparameter semantics, numerical-method order) are not
AST-local; syntactic mutator design goals (operator typos, off-by-one,
negation flips) have hit probability for domain errors much smaller than
systematic semantic-mutator triggering, and lack repeatability.
\textbf{Conclusion.} Syntactic tools occasionally producing mutants
satisfying (a)(b)(c) is stochastic byproduct, not repeatable, not
carrying engineering value. Systematic semantic mutation requires
(a)(b)(c) to be design intent.

\subsection{Per-class operator specialisations}\label{b.2-per-class-operator-specialisations}

{%
\begin{longtable}[]{@{}
  >{\raggedright\arraybackslash}p{(\linewidth - 2\tabcolsep) * \real{0.5000}}
  >{\raggedright\arraybackslash}p{(\linewidth - 2\tabcolsep) * \real{0.5000}}@{}}
\caption{Per-class operator specialisations.}\label{tab:p2-17}\\
\toprule\noalign{}
\begin{minipage}[b]{\linewidth}\raggedright
Operator
\end{minipage} & \begin{minipage}[b]{\linewidth}\raggedright
Representative specialisations across PUTs
\end{minipage} \\
\midrule\noalign{}
\endfirsthead
\endhead
\bottomrule\noalign{}
\endlastfoot
\textbf{mut\_C} Conservation-breaking & A1 Lorenz: add $\varepsilon_{\mathrm{drift}}$ to RHS
(slow Hamiltonian drift); A2 LU: decomposition omits k+1-th multiplier
(breaks det conservation); B1 Beta-Bin: posterior omits normalisation;
C1 GPR: covariance omits positive-definite diagonal term; D1 MLP:
backprop omits one gradient term. \\
\textbf{mut\_M} Monotonicity-breaking & A3 FDM: $\Delta$t coefficient
occasionally negative; B2 MCMC: acceptance min(1, r) $\to$ min(0.95, r); C2
PCE: high-order coefficient sort inserts inversion; D2 SVM:
decision-function sign flips near boundary. \\
\textbf{mut\_G} Convergence-breaking & A1 Lorenz: RK4 $\to$ 1.5-order
hybrid; A3 FDM: 2nd-order difference $\to$ 1st-order; B3 MC: doubling sample
size does not 1/N-reduce variance; C3 NN-Surr: training-epoch
truncation. \\
\textbf{mut\_T} Trajectory-distorting & A1 Lorenz: state-vector y/z
swap; B2 MCMC: insert independent-sampling segment; C3 NN-Surr:
training-target slow phase shift; D1 MLP: hidden-layer periodic-mask
activation. \\
\textbf{mut\_F} Fidelity-order-breaking & A2 LU: partial pivoting
degrades to no pivoting; C1 GPR: length-scale switches to coarse prior;
C2 PCE: high-order term randomly retains low-order; D3 LR:
regularisation occasionally large. \\
\end{longtable}
}

\textbf{Per-class HP / OS / TF substitution rules} (illustrative).
\textbf{HP} on class a: tolerance / max\_iter; on class c: GPR
\texttt{noise\_level} / \texttt{length\_scale}; on class d: MLP
\texttt{hidden\_dim} / dropout. \textbf{OS} on class a: numerical-linalg
API swaps (\texttt{det} $\leftrightarrow$ \texttt{sum(diag)}); on class b:
probability-distribution sampling swaps; on class c: surrogate-class
swaps (GPR $\leftrightarrow$ RBF $\leftrightarrow$ NN). \textbf{TF} on class a: integration-order
changes (RK4 $\to$ Euler); on class b: MC estimator changes.

\textbf{Operator-level cross-table with cosmic-ray defaults}. The 12
cosmic-ray default operators are AST-local:
\texttt{NumberReplacer} (Num/Constant, CE partial, $\bigtriangleup$ literal values
only); \texttt{ReplaceArithmeticOperator} (BinOp);
\texttt{ReplaceComparisonOperator} (Compare);
\texttt{ReplaceLogicalOperator} (BoolOp); \texttt{ReplaceUnaryOperator}
(UnaryOp); \texttt{ReplaceTrueFalse} (NameConstant, CE keyword, $\bigtriangleup$
Bool literal); \texttt{BreakContinueReplacer} (Break/Continue);
\texttt{RemoveDecorator}; \texttt{RemoveExceptHandler};
\texttt{ZeroIterationForLoop} (For); \texttt{ReplaceIfBlock} (If);
\texttt{MutateSubscript}. None correspond to OS / HP / TF / SI:

{%
\begin{longtable}[]{@{}
  >{\raggedright\arraybackslash}p{(\linewidth - 4\tabcolsep) * \real{0.3333}}
  >{\raggedright\arraybackslash}p{(\linewidth - 4\tabcolsep) * \real{0.3333}}
  >{\raggedright\arraybackslash}p{(\linewidth - 4\tabcolsep) * \real{0.3333}}@{}}
\caption{Per-class operator specialisations.}\label{tab:p2-18}\\
\toprule\noalign{}
\begin{minipage}[b]{\linewidth}\raggedright
Semantic class
\end{minipage} & \begin{minipage}[b]{\linewidth}\raggedright
Tool support
\end{minipage} & \begin{minipage}[b]{\linewidth}\raggedright
Coverage
\end{minipage} \\
\midrule\noalign{}
\endfirsthead
\endhead
\bottomrule\noalign{}
\endlastfoot
\textbf{OS} API replacement & None & $\bigtriangleup$ 88.33\% disjoint (incidental
low-complexity hits) \\
\textbf{HP} & None & $\times$ tool inexpressible (0/72) \\
\textbf{TF} numerical-method order & None & $\times$ tool inexpressible
(0/54) \\
\textbf{SI / CF} structural injection & None & $\times$ tool inexpressible
(0/33) \\
\end{longtable}
}

\textbf{Operator-level conclusion.} All 13 default classes remain
AST-local (BinOp / Compare / BoolOp / UnaryOp / NameConstant / Number /
Subscript / If / For / Break / Continue / decorator / except). Categorically, no
entry recognises sklearn / scipy hyperparameter semantics (HP),
numerical-method order (TF), or control-flow intent (SI / CF). For OS,
low-complexity sub-expressions can be incidentally hit by BinOp;
empirical 11.67\% (the main manuscript's AST-overlap audit) is consistent with 88.33\% AST-disjointness
lower bound.

\subsection{Syntactic-Mutation Pilot Study}\label{b.3-cosmic-ray-a1-single-put-pilot-pre-12-put-empirical}

Before the 12-PUT generalisation in the main manuscript's AST-overlap audit, a single-PUT lightweight
empirical used \path{scripts/run_cosmic_ray_a1.sh} on a1 (file
934 B, AST-bound). Outcome: mutants\_generated \textasciitilde tens; $\geq$
90\% belong to BinOp / Compare / 13 default classes; small
killed-by-\texttt{test\_a1.py} ratio (PUT unit tests are
output-shape/type sanity checks; most cosmic-ray mutants not killed), an
empirical manifestation of the semantic-mutation criteria in the main manuscript: syntactic-tool mutants
are not aligned with MR-violation detection. The 12-PUT generalisation
(the main manuscript's AST-overlap audit) supersedes this pilot but corroborates the same conclusion at
scale.

\subsection{Expected Root-Cause Risk Profile}\label{b.4-expected-lrca-risk-profile}

\textbf{PUT-class $\times$ LRCA-layer risk weights.} A numeric: C2 dominant
($\star\star\star$ on L1 tolerance). B probabilistic: C4 dominant ($\star\star\star$ on L3). C
surrogate: C3 dominant ($\star\star\star$ on L2 OOD; $\star\star$ tolerance). D ML: C3 + C4
mixed ($\star\star\star$ L2; $\star\star$ L3).

\textbf{Operator-PUT root-cause hotspots (expected).} A: mut\_C/G/F
dominantly C2 (numerical tolerance), mut\_M/T C1 (true semantic). B:
dominantly C4 (statistical-assumption violations). C: mut\_C/M
dominantly C2/C1 on c1/c2, mut\_T/F dominantly C3 on c1/c2; c3 (NN-Surr)
row dominantly C5 (artefact). D: d1 row dominantly C5 (training-noise
artefact); d2/d3 dominantly C1 / C3 mixed.

\textbf{Expected suspect\_share thresholds.} A: 0.10-0.20 (acceptance $\leq$
0.25); B: 0.20-0.35 ($\leq$ 0.40); C: 0.20-0.30 ($\leq$ 0.35); D: 0.25-0.40 ($\leq$
0.45). 60-cell average: $\leq$ 0.20 (= H4; acceptance $\leq$ 0.25).

\subsection{Engineering-significance mapping}\label{b.5-engineering-significance-mapping}

j = k aligned diagonal is the H2 threshold-test slice; j $\neq$ k
off-diagonal is control; 9 vacant cells ($\circ$, no MR exercised) are not formally
adjudicated and are repurposed in the main manuscript's R\_sem/R\_kill decoupling discussion as descriptive evidence for
R\_sem / R\_kill decoupling. The v3 $\to$ v4 micro-change
($\Delta\delta = -0.009$, paired-role bootstrap 95\% CI [-0.238, 0.207])
attributes to LLM-source diversity under a fixed prompt; this is the
source-axis null contrast underpinning the main manuscript's aligned-versus-cross results.

\subsection{Default-Operator Overlap between Syntactic Mutation Tools}\label{b.6-mutmut-vs-cosmic-ray-default-operator-overlap}

Reviewer-requested manual operator-class cross-reference between
cosmic-ray's 13 default operators (cosmic-ray 8.4.6,
\path{cosmic_ray.operators} package) and mutmut's 14 default
operators (mutmut current \texttt{main},
\path{src/mutmut/mutation/mutators.py}). Both default-operator sets
are first-order AST-local replacements; neither implements a
cross-function-boundary, domain-knowledge-dependent, or
algorithmic-class-changing mutation (the three necessary conditions in the semantic-mutation criteria of the main manuscript
for semantic mutation).

{\tiny\def\LTcaptype{table} % do not increment counter
\begin{longtable}[]{@{}
  >{\raggedright\arraybackslash}p{(\linewidth - 8\tabcolsep) * \real{0.1800}}
  >{\raggedright\arraybackslash}p{(\linewidth - 8\tabcolsep) * \real{0.2200}}
  >{\raggedright\arraybackslash}p{(\linewidth - 8\tabcolsep) * \real{0.2200}}
  >{\raggedright\arraybackslash}p{(\linewidth - 8\tabcolsep) * \real{0.1600}}
  >{\raggedright\arraybackslash}p{(\linewidth - 8\tabcolsep) * \real{0.2200}}@{}}
\caption{Mutmut vs cosmic-ray default-operator overlap.}\label{tab:p2-19}\\
\toprule\noalign{}
\begin{minipage}[b]{\linewidth}\raggedright
Operator class
\end{minipage} & \begin{minipage}[b]{\linewidth}\raggedright
cosmic-ray default
\end{minipage} & \begin{minipage}[b]{\linewidth}\raggedright
mutmut default
\end{minipage} & \begin{minipage}[b]{\linewidth}\raggedright
Reaches the semantic-mutation criteria (a/b/c)?
\end{minipage} & \begin{minipage}[b]{\linewidth}\raggedright
Semantic-class reachability
\end{minipage} \\
\midrule\noalign{}
\endfirsthead
\endhead
\bottomrule\noalign{}
\endlastfoot
Numeric literal $\pm$1 & \texttt{number\_replacer} &
\texttt{operator\_number} & (a) no, (b) no, (c) no & partial CE only
(e.g.~\texttt{\_RHO=28.0\ $\to$\ 27.5}) \\
Binary arithmetic swap (+, $-$, $\times$, $\div$) &
\texttt{binary\_operator\_replacement} & \texttt{operator\_swap\_op}
(binary subset) & (a) no, (b) no, (c) no & partial OS only
(incidental) \\
Comparison swap (\textless, $\leq$, \textgreater, $\geq$, ==, !=) &
\texttt{comparison\_operator\_replacement} & \texttt{operator\_swap\_op}
(compare subset) & (a) no, (b) no, (c) no & none \\
Boolean swap (and $\leftrightarrow$ or) & \texttt{boolean\_replacer} &
\texttt{operator\_swap\_op} (bool subset) & (a) no, (b) no, (c) no &
none \\
Unary remove / swap & \texttt{unary\_operator\_replacement} &
\texttt{operator\_remove\_unary\_ops}, \texttt{operator\_swap\_op} & (a)
no, (b) no, (c) no & none \\
Keyword swap (True/False, etc.) & \texttt{keyword\_replacer} &
\texttt{operator\_keywords}, \texttt{operator\_name} & (a) no, (b) no,
(c) no & none \\
break $\leftrightarrow$ continue / return & \texttt{break\_continue} &
\texttt{operator\_keywords} (subset) & (a) no, (b) no, (c) no & none \\
Exception class swap & \texttt{exception\_replacer} & (none) & (a) no,
(b) no, (c) no & none \\
Decorator removal & \texttt{remove\_decorator} & (none) & (a) no, (b)
no, (c) no & none \\
Variable-binding insert / replace & \texttt{variable\_inserter},
\texttt{variable\_replacer} & \texttt{operator\_arg\_removal},
\texttt{operator\_assignment} & (a) no, (b) no, (c) no & none \\
no\_op insertion & \texttt{no\_op} & (none) & (a) no, (b) no, (c) no &
none \\
Zero-iteration \texttt{for} loop & \texttt{zero\_iteration\_for\_loop} &
(none) & (a) no, (b) no, (c) no & none \\
String content tweak (XX prefix, case flip) & (none) &
\texttt{operator\_string} & (a) no, (b) no, (c) no & none \\
Lambda body $\to$ \texttt{None}/\texttt{0} & (none) &
\texttt{operator\_lambda} & (a) no, (b) no, (c) no & none \\
Dict kwarg name (XX prefix) & (none) &
\texttt{operator\_dict\_arguments} & (a) no, (b) no, (c) no & none \\
Sym./asym. string-method swap & (none) &
string-method symmetric swap; string-method asymmetric swap & (a) no,
(b) no, (c) no & none \\
Augmented $\to$ simple assignment & (none) &
\texttt{operator\_augmented\_assignment} & (a) no, (b) no, (c) no &
none \\
\texttt{match} case removal & (none) & \texttt{operator\_match} & (a)
no, (b) no, (c) no & none \\
\end{longtable}
}

\textbf{Aggregate.} Cosmic-ray $\cap$ mutmut: \textasciitilde9 of 13
cosmic-ray operators have a near-equivalent in mutmut (numeric, binary,
compare, boolean, unary, keyword, break-continue, variable, plus partial
overlap on string-prefix vs \texttt{XX}); 4 cosmic-ray operators are
unique (exception, decorator, no\_op, zero-iteration-for); 6 mutmut
operators are unique (string content, lambda, dict-kwarg, string-method
swap, aug-assignment, match-case). The union is approximately 21
distinct first-order AST-local operator classes. None of the 21 can:

\begin{itemize}

\item
  \begin{enumerate}
  
  \item
    Cross a function-call or module-import boundary (e.g.~replace
    \texttt{det(M)} with \texttt{sum(np.diag(M))} requires knowing the
    equivalence on diagonal matrices, outside any default-operator
    scope);
  \end{enumerate}
\item
  \begin{enumerate}
  \setcounter{enumi}{1}
  
  \item
    Depend on domain knowledge for legality (e.g.~a hyperparameter swap
    that knows \texttt{noise\_level=1e-4} is the load-bearing knob in a
    Gaussian Process Regression PUT requires PUT-class awareness);
  \end{enumerate}
\item
  \begin{enumerate}
  \setcounter{enumi}{2}
  
  \item
    Change the algorithmic class (e.g.~swap
    \texttt{polynomial\ $\to$\ spline\ basis} rewrites the surrogate's
    mathematical structure).
  \end{enumerate}
\end{itemize}

The mutmut / cosmic-ray null intersection on the main manuscript's AST-overlap audit's HP/SI/TF classes
(zero overlap) therefore reflects a \textbf{structural property of
first-order AST-local mutation} rather than a tool-specific limitation,
supporting the preventive-defence framing in the main manuscript.

% [thematic break removed]

\section{Experimental Procedure Details}\label{c.-experimental-procedure-details}

\subsection{Cross-Source Generation Protocol}\label{c.1-cross-source-v4-protocol-full-specification}

\textbf{Two-stage ablation.}

{%
\begin{longtable}[]{@{}
  >{\raggedright\arraybackslash}p{(\linewidth - 6\tabcolsep) * \real{0.2500}}
  >{\raggedright\arraybackslash}p{(\linewidth - 6\tabcolsep) * \real{0.2500}}
  >{\raggedright\arraybackslash}p{(\linewidth - 6\tabcolsep) * \real{0.2500}}
  >{\raggedright\arraybackslash}p{(\linewidth - 6\tabcolsep) * \real{0.2500}}@{}}
\caption{Cross-source generation protocol.}\label{tab:p2-20}\\
\toprule\noalign{}
\begin{minipage}[b]{\linewidth}\raggedright
Version
\end{minipage} & \begin{minipage}[b]{\linewidth}\raggedright
Mutant pool
\end{minipage} & \begin{minipage}[b]{\linewidth}\raggedright
c-class primary MR
\end{minipage} & \begin{minipage}[b]{\linewidth}\raggedright
Use
\end{minipage} \\
\midrule\noalign{}
\endfirsthead
\endhead
\bottomrule\noalign{}
\endlastfoot
v3 & Single-source (Claude Opus 4.6) & P-ord. (pre-registered) &
H2 baseline \\
v4 & \textbf{Cross-source} (Claude Opus 4.6 + GPT-5.4 + DeepSeek chat) &
P-ord. (held at pre-registered) & Isolate LLM-source diversity \\
\end{longtable}
}

\textbf{Protocol} (\path{scripts/cross_source_campaign.py}). (a) For
each (PUT, operator), Claude / GPT / DeepSeek run K = 3 trials with
identical prompt (the cross-source protocol in the main manuscript), temperature 0.7; \texttt{source\_tag}
propagated to filenames. (b) \textbf{V1-V4 mechanical validation}
(\path{src/p2/mutators/validation.py}): V1 syntax
(\texttt{ast.parse}), V2 executability, V3 non-triviality
(\texttt{\textbar{}y\_mutant\ -\ y\_original\textbar{}\ \textgreater{}\ 1e-6}),
V4 signature consistency. (c) \textbf{DeepSeek model:}
\texttt{deepseek-chat} (113 tokens/call, 1.8 s), not
\texttt{deepseek-v4-pro} (340 tokens/call, 11 s), quality-equivalent
on a2\_OS1 dry-run. (d) \textbf{Pool capacity:} 37 ops $\times$ 3 sources $\times$ 3
trials = 333 attempts; V1-V4 pass rate 89.5\%; 298 confirmed mutants, of
which 292 enter the v4 pool after final-stage rejection of duplicates and
QA failures. Sources contribute nearly equally: Claude 101 / GPT 98 / DeepSeek 99
(Phase A engineering finding). (e) \textbf{Sampling}
(\path{scripts/build_pools.py}, POOL\_VERSION = v4): per-PUT 30 max,
measured mean 24.3, range 10-30. c1 (GPR) has only 10 because c1\_HP1
/ c1\_CE1 have V1-V4 near-zero pass (WhiteKernel noise\_level 1e-4 $\to$
1e-1 perturbation has minimal output impact, triggers V3 non-trivial
failure, which supports the main manuscript's R\_sem/R\_kill decoupling discussion).

\textbf{Protocol asymmetry.} v4 does not invoke reviewer LLM (cost
/ speed priority; a follow-up study will rerun dual-blind on the v4
grid, \textasciitilde\$5-8 USD); v3 used the original Phase-1 dual-blind
(Claude gen + GPT-5.4 review + DeepSeek arbitration). A fraction of the
v3 $\to$ v4 source-axis shift may be slight v4 quality decline rather than
LLM-source diversity.

\subsection{Differential prompt protocol}\label{c.1.1-differential-prompt-protocol-future-work-commitment}

The v3 $\to$ v4 micro-change of -0.009 (paired-role bootstrap 95\% CI
[-0.238, 0.207]) may reflect ``three LLMs under
identical prompt'' rather than ``true upper bound of LLM diversity''.
Separation requires \emph{one tailored differential prompt per LLM}
(skeleton: \path{scripts/run_differential_prompt.py}).

\textbf{Design (2$\times$3 factorial, within-(PUT, operator)).} Factor A:
prompt template, 3 levels: \textbf{V\_canonical} (control, identical to
v4), \textbf{V\_persona} (per-LLM identity: Claude ``numerical-analysis
editor''; GPT ``scientific-software refactorer''; DeepSeek ``library-API
substitution specialist''), \textbf{V\_cot} (Claude
\texttt{\textless{}thinking\textgreater{}} tags; GPT ``step by step'';
DeepSeek stepped-reasoning). Factor B, LLM source (Claude / GPT-5.4 /
DeepSeek chat). K = 3 per cell; total 37 $\times$ 3 $\times$ 3 $\times$ 3 = \textbf{999
trials}.

\textbf{Exit criteria.} If max(delta) $-$ min(delta) \textless{} 0.05: the
v3 $\to$ v4 $-$0.009 source-axis null is robust. If $\geq$ 0.05 with prompt variant
pushing delta past 0.474: H2 revised from ``not met'' to
``prompt-conditional met''. If $\geq$ 0.05 but all delta \textless{} 0.474:
prompt sensitivity exists but H2 unchanged.

\textbf{Resources.} API \textasciitilde\$18-30 USD; wall 3-4 h at
concurrency 4; V1-V4 only, no reviewer LLM.

\subsection{Manual Pilot and Initial Prescreen}\label{c.2-manual-pilot-lineage-llm-client-configuration-and-l0-prescreen}

The original cross-source protocol used 60\% multi-LLM consensus (Claude Opus
4.6 + GPT-5.4 + DeepSeek chat unanimous) plus 40\% manual injection by
scientific-computing researchers. Each (mut\_j, S\_i) pair carried one
prompt template containing PUT source, semantic intent of mut\_j,
prohibition against revealing specific MRs (avoids prompt leakage), and
output requirements (syntactically correct, executable, single-point
modification, diff \textless{} 10 lines). Reproducibility parameters:
temperature 0.3 in manual-pilot stage (v4 cross-source uses 0.7 per
§C.1); fixed seed; 5 candidates per prompt with 2-3 retained.

\textbf{Double-blind review (Scheme C).} Generator LLM-G (Claude Opus)
and reviewer LLM-R (GPT-5.4) are cross-source with strict role
separation. LLM-R sees only PUT original + mutant code, unaware of
mut\_j category / MR content / generator identity, outputting a triple
(syntactic, executable, semantic-fault-injected). Double-confirmed
mutants enter pool; inconsistencies enter manual arbitration ($\leq$ 10\%).
From the double-confirmed pool, 20\% is sampled for manual review;
manual-vs-LLM-R inconsistency \textgreater{} 10\% triggers full manual
downgrade. Same-source bias mitigation stratifies LLM-R by PUT class.
Prompt-injection mitigation disables RAG / web search on the API path.

\textbf{L0 pool prescreen.} Each mutant passes static syntax (linters /
type checkers), a unit self-test (simple inputs produce finite output),
and double-blind sign-off. Target: 30-50 candidates per cell $\to$ 10-15
retained after prescreen.

\subsection{Three-Layer Root-Cause Diagnosis}\label{c.4-lrca-three-layer-execution-full}

Decision tree, multi-label priority (C5 \textgreater{} C4 \textgreater{}
C3 \textgreater{} C2 \textgreater{} C1), and output (\texttt{C1\_share},
\texttt{suspect\_share}) are stated in §A.2; the full pseudocode is
reproduced there. LRCA does \textbf{not} modify SMS; killed set is not
filtered. The main manuscript's LRCA-mass results use the best calibrated combination
(\texttt{ood\_band\ =\ 0.02,\ tolerance\_multiplier\ =\ 3.0,\ repeats\ =\ 20});
default-threshold results retained in \texttt{lrca\_60cell\_v3.json} as
control. Interface with the preventive-defence risk profile in the main manuscript: L0 prescan $\leftrightarrow$ §C.2 dual-LLM
double-blind; L1 tolerance $\leftrightarrow$ N = 20 repetition subprocess; L2 OOD $\leftrightarrow$
input-distribution definition; L3 statistical baseline $\leftrightarrow$ pre-check
protocol.

\subsection{Pilot calibration}\label{c.4.1-pilot-calibration-37-operators-k-1020}

A 2026 Q3 operator-level pilot ran 12 PUTs $\times$ 37 operators (36
primary PUT-class operator pairs plus one CF structural-injection
specialisation; the conceptual semantic family count remains five; 12
\texttt{is\_key=True} at K=20, 25 at K=10; 470 trials total). Pilot
precursor quantities: R\_sem (V1-V4 $\wedge$ operator\_match pass rate),
D\_impl (median pairwise AST + literal + identifier Jaccard distance
among confirmed mutants), R\_kill (proportion of confirmed mutants
killed by AVP under that PUT's primary MR). Aggregate (N=37): R\_sem
median 0.50 mean 0.468 (0/37 with R\_sem=0); D\_impl median 0.42 mean
0.392 (1/37 with D\_impl $\approx$ 0); R\_kill median 0.00 mean 0.189. All 37
operators produced $\geq$ 1 V1-V4 passing mutant, a necessary
precondition for H1; the pre-registered H1 threshold itself ($\geq$ 5
non-equivalent mutants on $\geq$ 9/12 PUTs for $\geq$ 4 of 5 operator
families) is nonetheless not met, as reported in the main text.

\textbf{Key R\_sem / R\_kill decoupling pattern.} R\_sem $\geq$ 0.9 $\wedge$ R\_kill
= 0 combinations all on HP / TF / OS operators in c / d classes (e.g.,
c1\_CE1 noise 1e-4$\to$1e-1 with R\_sem 1.00, R\_kill 0.00; c3\_HP1
relu$\to$tanh, c3\_TF1 max\_iter 1000$\to$5, d1\_HP1 MLP $\alpha$, d3\_HP1 LR C).
R\_sem-high $\wedge$ R\_kill = 1 combinations all on CE / OS operators in a / b
classes (a2\_CE1 LU det, a2\_OS1 prod$\to$sum, b1\_OS1 $\alpha$/$\beta$ swap, d1\_TF1
label flip). This is empirical evidence for R\_sem / R\_kill decoupling at cell level.

% [thematic break removed]

\section{Statistical Analysis Details}\label{d.-statistical-analysis-details}

\subsection{Semantic-Mutation Score Variants and Construction Lemmas}\label{d.1-sms-variants-and-construction-lemmas}

\textbf{SMS\_unfiltered (reviewer comparison metric):}

\[
\mathrm{SMS}^{\mathrm{unfiltered}}_{i,k,j} := \frac{|\mathrm{killed}_{i,k,j}|}{|\mathrm{mut}_j(S_i)| - |\mathrm{equiv}_{i,k,j}|}
\]

(Same form as the primary metric; here \texttt{killed} does not distinguish C1 from C2--C5.)

The appendix of the reproducibility package provides cell-by-cell
difference between SMS\_unfiltered and SMS; if relative difference
\textless{} 5\%, this confirms LRCA does not affect robustness of
primary conclusions.

\textbf{Primary-table construction.}

{%
\begin{longtable}[]{@{}
  >{\raggedright\arraybackslash}p{(\linewidth - 4\tabcolsep) * \real{0.3333}}
  >{\raggedright\arraybackslash}p{(\linewidth - 4\tabcolsep) * \real{0.3333}}
  >{\raggedright\arraybackslash}p{(\linewidth - 4\tabcolsep) * \real{0.3333}}@{}}
\caption{SMS variants and construction lemmas.}\label{tab:p2-21}\\
\toprule\noalign{}
\begin{minipage}[b]{\linewidth}\raggedright
RQ
\end{minipage} & \begin{minipage}[b]{\linewidth}\raggedright
Primary statistics
\end{minipage} & \begin{minipage}[b]{\linewidth}\raggedright
Reporting format
\end{minipage} \\
\midrule\noalign{}
\endfirsthead
\endhead
\bottomrule\noalign{}
\endlastfoot
RQ3 & inst\_rate, equiv\_rate, C1\_share, survive\_rate & 60-cell
heatmap + 4-class marginals \\
RQ4 & aligned-SMS vs cross-SMS; sparse $\circ$ vs dense \textbullet\textbullet equiv\_rate &
Cliff's delta + odds ratio + 95\% bootstrap CI \\
RQ5 & $\Delta$SMS\_c (c $\in$ \{A,B,C,D\}); CV($\Delta$SMS) & sign test (df=3) +
descriptive forest plot \\
RQ5 (desc.) & Spearman $\rho$ + Kendall $\tau$ for SMS vs PC & scatter + dual-metric
ranking comparison \\
\end{longtable}
}

\textbf{Multiple comparisons.} No cell-level family of p-values is
claimed across the 60 cells (main text, statistical pipeline); the
per-class Friedman tests carry a Bonferroni $\times$ 4 correction.
N = 20 bootstrap intervals:
1,000-iteration 95\% CI (B = 10,000 for the headline H2 delta CI).

\subsection{Effect-Size Sensitivity and Transformation Invariance}\label{d.2-cliffs-delta-cutoff-sensitivity-and-logit-transformation-invariance}

\textbf{H4 cutoff sensitivity.} Dense grid scan (cutoff
$\in$ \{0.05, 0.10, \ldots, 0.50\}, step 0.05) on v4 data
(\path{scripts/h5_sensitivity.py},
\path{data/results/h5_sensitivity_v4.json}):

{%
\begin{longtable}[]{@{}lll@{}}
\caption{Effect-size sensitivity and logit-transformation invariance.}\label{tab:p2-22}\\
\toprule\noalign{}
cutoff & h5\_cells\_pass & h5\_pass\_ratio \\
\midrule\noalign{}
\endfirsthead
\endhead
\bottomrule\noalign{}
\endlastfoot
0.05 & 12 / 60 & 20.0\% \\
0.10 & 12 / 60 & 20.0\% \\
0.15 & 12 / 60 & 20.0\% \\
\textbf{0.20 (paper)} & \textbf{12 / 60} & \textbf{20.0\%} \\
0.25 & 12 / 60 & 20.0\% \\
0.30 & 12 / 60 & 20.0\% \\
0.35 & 12 / 60 & 20.0\% \\
0.40 & 12 / 60 & 20.0\% \\
0.45 & 13 / 60 & 21.7\% \\
0.50 & 13 / 60 & 21.7\% \\
\end{longtable}
}

\textbf{Conclusion.} H4 verdict (not met) is intrinsic data property,
independent of cutoff choice. v4 suspect\_share distribution is severely
bimodal: median 1.0, mean 0.79; 48 cross cells at suspect\_share $\approx$ 1, 12
aligned cells in low end; the {[}0.20, 0.80{]} interior is nearly empty.
Specific value 0.20 is \textbf{not load-bearing}.

\textbf{Logit-transformation invariance.} Cliff's delta is rank-based
(function of $U / (n_1 \cdot n_2)$), mathematically invariant under any strictly
monotone transformation (Romano 2006). Logit is strictly monotone on (0,
1), so $\delta_{\mathrm{logit}} \equiv \delta_{\mathrm{raw}}$ is a construction result. \textbf{Not
additional robustness evidence}; H2 robustness threats come from the main manuscript's 60-cell distribution results
zero-mass dominance, not metric-scale choice.

\textbf{Vacant-cell sensitivity.} Excluding the 9 vacant cells (cells
with no exercised MR) leaves the aligned-versus-cross verdict unchanged.
In v4 the analysis becomes n = 11 aligned and n = 40 cross with
$\delta = 0.323$ (95\% CI [-0.011, 0.648]); the matching v3 slice gives
$\delta = 0.270$ (95\% CI [-0.043, 0.595]). The exclusion widens the
interval because the design is smaller, but it does not move H2 across
the pre-registered Romano medium threshold.

\textbf{Primary-MR withdrawal permutation-null.} The auxiliary c-class
check pools 15 SMS values, shuffles labels over the $3 \times 5$ (PUT,
MR) grid, and applies the max-over-five withdrawal rule with 10,000
permutations (seed 42). The observed c-class mean is 0.314 versus null
mean 0.349, one-sided $p = 0.989$
(\path{data/results/c_class_permutation_v4.json}). This supports treating
the primary-MR convention as a protocol choice rather than as evidence of
a hidden c-class inflation effect.

\subsection{Power Analysis under Observed and Stipulated Effects}\label{d.3-power-analysis-plug-in-and-stipulated-alternative}

\textbf{Plug-in bootstrap.} With-replacement sample from observed
aligned (n=12) and cross (n=48) v4 SMS pools; N\_sim = 5,000, seed = 42
(\path{scripts/compute_rq2_power.py}, \path{rq2_power_v4.json}):

{%
\begin{longtable}[]{@{}lll@{}}
\caption{Power analysis under observed and stipulated effects.}\label{tab:p2-23}\\
\toprule\noalign{}
Threshold & Interpretation & Power \\
\midrule\noalign{}
\endfirsthead
\endhead
\bottomrule\noalign{}
\endlastfoot
delta \textgreater{} 0.000 & Any-effect & \textbf{0.997} \\
delta \textgreater{} 0.147 & Small & 0.966 \\
delta \textgreater{} 0.330 & Medium & 0.759 \\
\textbf{delta \textgreater{} 0.474} & \textbf{Large (H2)} &
\textbf{0.423} \\
\end{longtable}
}

Sample-size sweep ($n_{\mathrm{aligned}} \in \{6, 12, \ldots, 60\}$, $n_{\mathrm{cross}} = 4 \times n_{\mathrm{aligned}}$):
power for delta \textgreater{} 0 reaches 0.974 at n=6, 0.996 at 12,
plateaus thereafter. RQ4 primary analysis has very low sample-size
requirement; the bottleneck is the effect-size boundary itself, not
sampling noise. Even at $\delta_{\mathrm{truth}} \approx 0.474$, sample has only
\textasciitilde42\% chance of detecting under plug-in; but this does
\textbf{not} mean ``insufficient power is the cause of H2 not being
met'', observed $\delta_{\mathrm{v3}} = 0.323 < 0.474$, observed
$\delta_{\mathrm{v4}} = 0.314 < 0.474$; increasing sample only narrows CI.

\textbf{Stipulated-alternative.} Implemented in
mixture-weight (\path{scripts/compute_rq2_power_stipulated.py},
N\_sim = 2,000): SMS has heavy ties at zero (45/60 = 0), raw shifting is
discontinuous (any $\varepsilon$ \textgreater{} 0 jumps delta from 0.314 to 0.74).
Mixture: aligned' = (probability w) sample from (observed\_aligned +
0.001) + (probability 1$-$w) sample from observed\_aligned, calibrate w so
E{[}delta{]} $\approx$ 0.474. Calibrated w = 0.094, realised E{[}delta{]} =
0.4746.

{%
\begin{longtable}[]{@{}lll@{}}
\caption{Power analysis under observed and stipulated effects.}\label{tab:p2-24}\\
\toprule\noalign{}
Stipulated truth & Criterion & Power \\
\midrule\noalign{}
\endfirsthead
\endhead
\bottomrule\noalign{}
\endlastfoot
0.474 & $\hat{\delta} \geq 0.474$ (point-estimate) & \textbf{0.491} \\
0.474 & 95\% CI lower \textgreater{} 0 (any-effect) & 0.868 \\
\end{longtable}
}

Even at $\delta_{\mathrm{truth}} = 0.474$, the (12, 48) design returns $\hat{\delta} \geq
0.474$ in only \textasciitilde49\% of replications. This
\emph{supports} the main manuscript's aligned-versus-cross framing that ``not met'' is point-estimate, not
effect-size.

\subsection{Per-Class Friedman Tests}\label{d.4-friedman-per-class-discussion-of-small-n-concordance}

Per-class Friedman (3 SUTs $\times$ 5 MRs) with Bonferroni $\times$ 4: a chi\^{}2=4.00
raw 0.406 adj 1.000, W=0.333; b chi\^{}2=10.34 raw \textbf{0.035} adj
\textbf{0.140}, W=\textbf{0.862}; c chi\^{}2=5.60 raw 0.231 adj 0.924,
W=0.467; d chi\^{}2=5.00 raw 0.287 adj 1.000, W=0.417. After correction,
no per-class result remains significant at family-wise alpha = 0.05.
Even b-class W = 0.862 (Cohen 1988 large concordance) is insufficient at
N = 3 per class. \textbf{H3 verdict rests on the main manuscript's cross-class consistency sign test}; per-class
Friedman is sensitivity / descriptive only. Friedman main effect on full
60 cells (chi\^{}2 = 16.76, p = 0.0022) confirms MR rank differences but
is logically independent from H3 cross-class consistency.

\textbf{Mixed-effects unavailability.} Primary
\texttt{sms\ \textasciitilde{}\ C(class)\ +\ C(operator)\ +\ C(class):C(operator)\ +\ (1\ \textbar{}\ put)}
returned Singular (insufficient column rank for class $\times$ operator
interaction; N=60 too small for 11-d fixed-effects + 12 PUT random
intercepts). Fallback
\texttt{sms\ \textasciitilde{}\ C(class)\ +\ C(operator)\ +\ (1\ \textbar{}\ put)}
had PUT random-intercept variance hit boundary 0 (degenerates to OLS).
Fallback class p-values (descriptive only): b vs a 0.275 / c vs a 0.892
/ d vs a 0.991. RQ5 conclusion shifts to four-piece presentation (class
means + sign test + Friedman + forest plot).

\subsection{Pattern-coverage operationalisation}\label{d.5-pattern-coverage-operationalisation}

Per SUT, (MR\_k, R\_outcome $\in$ \{True, False\}) binary-tuple coverage: 5
MRs $\times$ 2 outcomes = 10 cells, PC = \#triggered / 10. Simplest baseline
(per-SUT granularity, not distinguishing mutants). PC range over 12
SUTs: {[}0.500, 1.000{]}, mean 0.733 (\texttt{paper\_numbers\_v4.json}).
Pairing per PUT: Spearman rho = 0.163; Kendall tau = 0.136 (descriptive
only; this row carries no inferential licence in the main manuscript's
inference-permissions table, so no $p$-value is attached).

\textbf{Power caveat.} At n = 12, Spearman 95\% CI $\approx$ {[}-0.5, +0.6{]}
(Fisher-z approximation); cannot distinguish ``zero correlation'' from ``moderate positive'' or
``moderate negative''. The descriptive rho / tau mean ``no correlation detected
at n = 12'', not ``correlation does not exist''.

\textbf{Within-class descriptive.} b2 PC = 1.0 with mean SMS = 0.067; b1
PC = 0.7 with mean SMS = 0.20. c3 PC = 1.0 with mean SMS = 0.14; c1 / c2
PC = 0.7-0.8 with mean SMS = 0.0. Within the same class, higher-PC PUT
has lower SMS, contradicts naive ``more PC kills more mutants''
assumption. This pattern is consistent with orthogonality, but n = 12
cannot confirm it. A confirmatory study would need n $\geq$ 30 and a PC
definition that incorporates the mutant dimension.

% [thematic break removed]

\section{Stakeholder Deployment Considerations}\label{e.-stakeholder-deployment-considerations-single-output-kernels-scope}

\subsection{Air-Gap Deployment Constraints}\label{e.1-air-gap-incompatibility-full-justification}

The stakeholder-analysis workflow in the main manuscript depends on external LLM API calls (Claude / GPT /
DeepSeek), incompatible with most regulated air-gapped V\&V workflows:

{%
\begin{longtable}[]{@{}
  >{\raggedright\arraybackslash}p{(\linewidth - 4\tabcolsep) * \real{0.3333}}
  >{\raggedright\arraybackslash}p{(\linewidth - 4\tabcolsep) * \real{0.3333}}
  >{\raggedright\arraybackslash}p{(\linewidth - 4\tabcolsep) * \real{0.3333}}@{}}
\caption{Air-gap deployment constraints.}\label{tab:p2-25}\\
\toprule\noalign{}
\begin{minipage}[b]{\linewidth}\raggedright
Domain
\end{minipage} & \begin{minipage}[b]{\linewidth}\raggedright
Standard
\end{minipage} & \begin{minipage}[b]{\linewidth}\raggedright
Air-gap requirement
\end{minipage} \\
\midrule\noalign{}
\endfirsthead
\endhead
\bottomrule\noalign{}
\endlastfoot
Nuclear & IEC 60880 & Air-gapped build for safety-critical software \\
Aerospace & DO-178C & Tool-qualified offline build for DAL-A/B; LLM API
outside qualified-tool envelope \\
Medical & IEC 62304 & Class C software lifecycle requires offline
reproducible build \\
Automotive & ISO 26262 & ASIL-D requires offline tool chain; cloud LLM
non-compliant \\
\end{longtable}
}

\textbf{Mitigation paths.} (a) Self-hosted open-weight LLMs (Llama
/ DeepSeek local inference), capability gap vs API-served frontier
models; quantitative impact on SMS untested. (b) Offline-cached
pre-generated mutant pools with commit-hash + raw-response signature
locking, adopters reproduce a published pool offline (this paper's
package supports this), but generating new pools requires external LLM
access. SMS for air-gapped industrial deployment is a future research
direction.

\subsection{Quarterly batch audit workflow}\label{e.2-quarterly-batch-audit-workflow}

The original per-PR GitHub Actions YAML was removed because its PR-CI
threshold 0.10 contradicted the Appendix~E.3 audit threshold 0.20 and propagated
stakeholder-facing CI assumptions into adopters' pipelines.

\textbf{Replacement: quarterly batch audit:}

{%
\begin{longtable}[]{@{}
  >{\raggedright\arraybackslash}p{(\linewidth - 4\tabcolsep) * \real{0.3333}}
  >{\raggedright\arraybackslash}p{(\linewidth - 4\tabcolsep) * \real{0.3333}}
  >{\raggedright\arraybackslash}p{(\linewidth - 4\tabcolsep) * \real{0.3333}}@{}}
\caption{Quarterly batch audit workflow.}\label{tab:p2-26}\\
\toprule\noalign{}
\begin{minipage}[b]{\linewidth}\raggedright
Step
\end{minipage} & \begin{minipage}[b]{\linewidth}\raggedright
Description
\end{minipage} & \begin{minipage}[b]{\linewidth}\raggedright
Cost
\end{minipage} \\
\midrule\noalign{}
\endfirsthead
\endhead
\bottomrule\noalign{}
\endlastfoot
1 & Offline mutant-pool generation (per PUT 24-30 $\times$ 12 PUTs $\approx$ 300) & LLM
API \$5-15; \textasciitilde30 min \\
2 & \texttt{sms\_campaign.py\ -\/-track\ 2} & 4-core laptop
\textasciitilde10-20 min \\
3 & \texttt{run\_lrca.py} & \textless{} 1 min \\
4 & MR-design team review of MRs with aligned-cell SMS significantly
below historical median (0.275) & 1-2 h meeting \\
\end{longtable}
}

Quarterly: API \textasciitilde\$10 + automation 30 min + manual 1-2 h $\approx$
0.5 person-day. Affordable. Per-PR gating not recommended (LLM latency
5-30 s + cost + non-determinism + air-gap incompatibility); adopters
needing per-PR should use coverage / unit-test gates. SMS does not
replace domain-expert MR physical-reasonableness judgement nor
product-requirements review.

\subsection{Verification-and-Validation Documentation}\label{e.3-vv-documentation-conceptual-complementarity-with-asme-vv-20-2009}

This appendix subsection makes \textbf{no normative claim} toward NRC,
FDA, or ISO 26262 review teams. Within single-output kernels scope, no
traceable mapping exists to current IEC 60880 / ISO 26262 / DO-178C /
ASME V\&V 20-2009 normative bodies.

V\&V documentation for scientific computing software (per ASME V\&V
20-2009 §3 and similar guides) requires quantifiable test-adequacy
evidence; current documentation often relies on code coverage + MR lists
+ SME signatures. SMS may appear as \textbf{research-grade supplementary
evidence} alongside coverage and MR lists: (a) aligned-cell SMS per
critical PUT; (b) LRCA three-layer diagnosis (C1 readout / C2-C5
attribution); (c) 60-cell matrix visualisation. Reviewers can
independently run \texttt{REPRODUCIBILITY.md}.

Original threshold recommendations (aligned-cell SMS $\geq$ 0.20 / 0.30 +
C1\_share $\leq$ 0.20) were removed: no normative backing in IEC / ISO /
ASME; misreads as enforce-ready. \textbf{SMS reported as descriptive
supplementary evidence}, with adequacy judgements left to V\&V reviewers
case-by-case.

Conceptually complementary to ASME V\&V 20-2009 §3 code verification
(numerical-solver correctness) vs SMS (MR-set fault-detection adequacy).
Substantial scale gap to multi-module CFD codes; incorporation into V\&V
standards body would require large-scale empirics + multi-year
dialogue with ASME V\&V 20 / IEEE 1012 / IEC 60880 committees.
\textbf{No advocacy that SMS enter any normative certification system in
2027.} SMS does not replace SME signatures, FMECA, PIRT, system-level
V\&V, or ASME V\&V 20 §3 numerical-solver verification, it is one
link in the evidence chain, not a single-point qualification.

% [thematic break removed]

\section{Threats to Validity: Detailed Mitigation}\label{f.-threats-to-validity-detailed-mitigation}

\subsection{Internal threats}\label{f.1-internal-threats}

\textbf{Hoeffding-style false-equivalence bound (referenced by
main-text Limitation 1).} For a mutant whose true per-sample
disagreement probability against the original program under $D_S$ is
at least $p_0$ (at the E2 tolerance), the probability that all
$K_{\mathrm{eq}}$ independent E2 samples agree, producing a false
equivalence verdict, is at most $(1-p_0)^{K_{\mathrm{eq}}} \leq
\exp(-K_{\mathrm{eq}}\,p_0)$. At $K_{\mathrm{eq}} = 1000$: $p_0 =
0.005$ gives $\leq e^{-5} \approx 0.7\%$; $p_0 = 0.01$ gives $\leq
e^{-10} \approx 4.5 \times 10^{-5}$. The bound conditions on $p_0$ and
says nothing about mutants whose disagreement is concentrated below the
E2 tolerance; those are governed by the E1 arm and remain the residual
risk the K\_eq sweep (deferred) would quantify.

{%
\begin{longtable}[]{@{}
  >{\raggedright\arraybackslash}p{(\linewidth - 6\tabcolsep) * \real{0.2500}}
  >{\raggedright\arraybackslash}p{(\linewidth - 6\tabcolsep) * \real{0.2500}}
  >{\raggedright\arraybackslash}p{(\linewidth - 6\tabcolsep) * \real{0.2500}}
  >{\raggedright\arraybackslash}p{(\linewidth - 6\tabcolsep) * \real{0.2500}}@{}}
\caption{Internal threats.}\label{tab:p2-27}\\
\toprule\noalign{}
\begin{minipage}[b]{\linewidth}\raggedright
Threat area
\end{minipage} & \begin{minipage}[b]{\linewidth}\raggedright
Threat
\end{minipage} & \begin{minipage}[b]{\linewidth}\raggedright
Mitigation
\end{minipage} & \begin{minipage}[b]{\linewidth}\raggedright
Residual
\end{minipage} \\
\midrule\noalign{}
\endfirsthead
\endhead
\bottomrule\noalign{}
\endlastfoot
LLM generation & LLM generation reproducibility & §C.2 reproducibility parameters in
package; dual-LLM cross-source + 20\% manual; reproducibility rests on
archived raw responses and frozen mutant pools (no user-facing seed) & LLM training data may be updated
post-submission \\
Equivalence approximation & Probabilistic equiv (K\_eq=1000) & the main manuscript's equivalence-judgement section declares probabilistic
approximation; Hoeffding-style false-equiv bound & K\_eq sweep $\in$
\{500,1000,2000\} deferred to future work \\
Shared infrastructure & Shared-infrastructure dependency & AVP version pinned to the upstream commit hash;
package embeds AVP source; the arXiv infrastructure technical report is disclosed in the main text &
None within scope \\
Root-cause diagnosis & LRCA multi-label boundary & §A.2 priority
C5 \textgreater{} C4 \textgreater{} C3 \textgreater{} C2 \textgreater{} C1;
multi-label co-occurrence table in package; root\_cause is ``likely root
cause'' not definitive & C2-C5 may multi-trigger on B/D classes \\
Pool size & Mutant-pool size & Pool expanded to mean 17.4: delta moved 0.321 $\to$
0.323, CI narrowed; \textbf{effect size is intrinsic ceiling, not pool
dilution} & Some PUTs \textless{} 30 (cache-bound) \\
LLM non-determinism & LLM non-determinism & Multi-turn de-dup; K=10/20 repetitions; raw
prompt+response committed for direct reuse & Claude Opus subscription
lacks seed control \\
\textbf{Protocol asymmetry} & \textbf{Protocol-implementation gap v3 vs v4} & v4 used 3
providers (Claude 101 / GPT 98 / DeepSeek 99) without dual-blind; v4
C1\_share 0.209 \textgreater{} v3 0.164 weakly opposes ``v4 quality
decline'' & Complete separation requires a follow-up dual-blind rerun on
the v4 grid \\
\end{longtable}
}

The protocol-asymmetry threat (dual-blind vs V1-V4 only) applies to the
v3 $\to$ v4 source-axis contrast ($\Delta\delta = -0.009$, paired-role
bootstrap 95\% CI [-0.238, 0.207]).

\subsection{External / construct / conclusion threats}\label{f.2-external-construct-conclusion-threats}

{%
\begin{longtable}[]{@{}
  >{\raggedright\arraybackslash}p{(\linewidth - 6\tabcolsep) * \real{0.2500}}
  >{\raggedright\arraybackslash}p{(\linewidth - 6\tabcolsep) * \real{0.2500}}
  >{\raggedright\arraybackslash}p{(\linewidth - 6\tabcolsep) * \real{0.2500}}
  >{\raggedright\arraybackslash}p{(\linewidth - 6\tabcolsep) * \real{0.2500}}@{}}
\caption{External / construct / conclusion threats.}\label{tab:p2-28}\\
\toprule\noalign{}
\begin{minipage}[b]{\linewidth}\raggedright
Threat area
\end{minipage} & \begin{minipage}[b]{\linewidth}\raggedright
Threat
\end{minipage} & \begin{minipage}[b]{\linewidth}\raggedright
Class
\end{minipage} & \begin{minipage}[b]{\linewidth}\raggedright
Mitigation
\end{minipage} \\
\midrule\noalign{}
\endfirsthead
\endhead
\bottomrule\noalign{}
\endlastfoot
Representativeness & 12-PUT representativeness & External & follows the shared site selection;
open class principle (cls extensible); scaling to MD/QC/CFD \\
Statistical power & Cross-class statistical power & Conclusion & H3 framed exploratory;
mixed-effects unavailable (Singular at N=60); shifted to class means +
sign test + Friedman + forest plot \\
LLM homogeneity & LLM homogeneity bias & Construct & §C.2 three-LLM rotation +
PUT-class subdivision + 20\% manual sampling; third-party LLM family
validation deferred \\
LLM-source shift & LLM-source distributional shift & External & Hit distribution
DeepSeek 11/15, Claude 4/15, GPT 0/15; the main manuscript's AST-overlap audit argument is categorical
AST-locality, not hit frequency; HP/SI/TF zero-overlap across all 3
sources unaffected \\
Higher-order mutation & HOM equivalence & External & Tool-unreachability claim confined to
first-order syntactic tools; HOM empirical testing listed as future
work \\
\end{longtable}
}

Conclusion threats: multiple comparisons handled per the main-text
policy (per-class Bonferroni $\times$ 4; no cell-level family claimed;
§D.1); N=20 stability handled by 1,000-iteration
bootstrap CI (§D.3). Construct threat, SMS measures epistemological
semantic detection, not production engineering value (out of scope here).

% [thematic break removed]

\section{Degeneration to Classical Mutation Score: Full Proof}\label{g.-sms-ms-degeneration-theorem-full-proof}

\subsection{Notation cross-reference}\label{g.1-notation-cross-reference-with-2.1}

PUT $S_i$, mutation operator family mut (syntactic or semantic),
mutant set $\mathrm{mut}(S_i)$, MR-label set, equivalence tolerance $\varepsilon_{\mathrm{eq}}$, equivalence
sample size $K_{\mathrm{eq}}$, AVP tolerance $\varepsilon_{\mathrm{AVP}}$. Three-state decomposition:
$\mathrm{mut}(S)=\mathrm{killed} \cup \mathrm{equiv} \cup \mathrm{survive}$ (disjoint). SMS
formula:

\[
\mathrm{SMS}_{i,k,j} = \frac{|\mathrm{killed}_{i,k,j}|}{|\mathrm{mut}_j(S_i)| - |\mathrm{equiv}_{i,k,j}|}
\]

\subsection{Degenerate-limit definition}\label{g.2-degenerate-limit-definition-r-8-p1-3-revision-3-joint-conditions}

The degenerate limit $L = L_{\mathrm{equiv}} \wedge L_{\mathrm{killed}} \wedge L_{\mathrm{mut}}$
consists of three \textbf{joint conditions} (paired axes), each
controlling one layer of the SMS formula. L1-L6 are not 6 independent
axes; the pairing reflects dependency structure across the SMS formula
layers.

{%
\begin{longtable}[]{@{}
  >{\raggedright\arraybackslash}p{(\linewidth - 4\tabcolsep) * \real{0.3333}}
  >{\raggedright\arraybackslash}p{(\linewidth - 4\tabcolsep) * \real{0.3333}}
  >{\raggedright\arraybackslash}p{(\linewidth - 4\tabcolsep) * \real{0.3333}}@{}}
\caption{Degenerate-limit definition.}\label{tab:p2-29}\\
\toprule\noalign{}
\begin{minipage}[b]{\linewidth}\raggedright
Joint condition
\end{minipage} & \begin{minipage}[b]{\linewidth}\raggedright
Component axes
\end{minipage} & \begin{minipage}[b]{\linewidth}\raggedright
Pairing rationale
\end{minipage} \\
\midrule\noalign{}
\endfirsthead
\endhead
\bottomrule\noalign{}
\endlastfoot
\textbf{$L_{\mathrm{equiv}}$} (Lemma G.1) & L1: $\varepsilon_{\mathrm{eq}} \to 0$; L2: $K_{\mathrm{eq}} \to \infty$ & L1 alone:
equiv stays probabilistic ($K_{\mathrm{eq}}$ cannot cover $D_S$). L2 alone: bitwise
equality diluted by $\varepsilon_{\mathrm{eq}}$. Both simultaneously degenerate equiv to
classical behavioural equivalence outside a D\_S-measure-zero set. \\
\textbf{$L_{\mathrm{killed}}$} (Lemma G.2) & L3: $\varepsilon_{\mathrm{AVP}}^k \to 0$; L4: MR-label set =
$\{\mathrm{MR}_{\mathrm{eq}}\}$, R(y,y') $\equiv$ y = y' & L3 alone: $\varepsilon_{\mathrm{AVP}} \to 0$ still allows
non-trivial MR (monotonicity, convergence). L4 alone: equality still
carries $\varepsilon_{\mathrm{AVP}}$ tolerance. Both simultaneously degenerate killed to
classical difference detection. \\
\textbf{$L_{\mathrm{mut}}$} (Lemma G.3) & L5: $\mathrm{mut}_j \to$ rule-based syntactic (Mothra
AOR/ROR/SDL/CRP); L6: cls(I) $\subseteq$ \{imperative deterministic\} & L5 alone:
syntactic ops on probabilistic/ML may still trigger domain-semantic
subsets. L6 alone: imperative programs still mutable by OS/HP/TF/SI.
Both simultaneously degenerate mut(S) to Jia \& Harman syntactic mutant
set. \\
\end{longtable}
}

\subsection{Lemmas for Three-State Decomposition}\label{g.3-lemmas-three-state-decomposition-degenerates-under-l}

\textbf{Lemma G.1 (equiv degeneration; measure-zero
qualification).} Under $L_{\mathrm{equiv}}$ (L1 $\wedge$ L2), semantic-class equivalence
(E1 $\wedge$ E2) degenerates to classical behavioural equivalence
\textbf{almost everywhere} w.r.t. measure $D_S$.

\emph{Proof.} E1 (type consistency) holds trivially under $\varepsilon_{\mathrm{eq}} \to 0$ (L6
makes imperative output spaces scalar/vector with static types). E2
($|S_i(x) - s'(x)| < \varepsilon_{\mathrm{eq}}$ for $K_{\mathrm{eq}}$ samples)
under L1 ($\varepsilon_{\mathrm{eq}} \to 0$) $\wedge$ L2 ($K_{\mathrm{eq}} \to \infty$ with measure-equivalent sampling)
is almost-everywhere equivalent to $\forall x \in D_S \setminus N: S_i(x) = s'(x)$, where
$N$ is a $D_S$-measure-zero set (continuous $D_S$: floating-point
pathological points / NaN propagation; discrete $D_S$: $N = \emptyset$). This
matches \citet{jia2011analysis} §3 classical equivalent-mutant definition
under measure-zero equivalence classes.

\textbf{Lemma G.2 (killed degeneration).} Under L3 $\wedge$ L4, killed
degenerates to classical difference detection.

\emph{Proof.} L4 restricts the MR-label set to \{MR\_eq\} with R(y, y') $\equiv$ y = y'. For
mr = (r, R) and mutant s', the $\mathrm{MR}_{\mathrm{eq}}$ violation condition $\exists x: S_i(x) \neq
s'(r(x))$ under L3 ($\varepsilon_{\mathrm{AVP}} \to 0$) becomes exact inequality. With r = id,
violation is $S_i(x) \neq s'(x)$, classical difference detection. With r
$\neq$ id, $\mathrm{MR}_{\mathrm{eq}}$ still requires $S_i(x) = s'(r(x))$ as reference oracle from
original program; no new state classifications introduced.

\textbf{Lemma G.3 (mut degeneration).} Under L5 $\wedge$ L6, mut\_j(S\_i)
degenerates to the syntactic mutant set of \citet{jia2011analysis}.

\emph{Proof.} L5 switches $\mathrm{mut}_j$ to rule-based syntactic operators (AOR,
ROR, SDL, CRP, UOI, standard Mothra/Proteum sets); L6 restricts PUTs
to imperative deterministic programs, excluding triggering conditions
for semantic operators on probabilistic/ML programs. Under this
configuration, $\mathrm{mut}_j(S_i)$ is the literature-defined syntactic mutant
set, independent of domain semantics.

\subsection{Degeneration Theorem: Detailed Proof}\label{g.4-theorem-9.1-detailed-proof}

\textbf{Theorem G (SMS $\to$ MS degeneration; stated in the main text
as the degeneration theorem).} In the degenerate limit
$L = L_{\mathrm{equiv}} \wedge L_{\mathrm{killed}} \wedge L_{\mathrm{mut}}$, almost everywhere
w.r.t. $D_S$,

\[ \mathrm{SMS}_{i,k,j} \xrightarrow{L} \mathrm{MS}_{i,j} := \frac{|\mathrm{killed}_{i,j}^{\mathrm{classic}}|}{|\mathrm{mut}_j^{\mathrm{syntax}}(S_i)| - |\mathrm{equiv}_{i,j}^{\mathrm{classic}}|} \]

\emph{Proof.} Combine Lemmas 9.1-9.3:

\begin{itemize}

\item
  Numerator: under L3 $\wedge$ L4, $\mathrm{killed}_{i,k,j} \to
  \mathrm{killed}_{i,j}^{\mathrm{classic}}$ (Lemma G.2); L4 makes $\mathrm{MR}_{i,k}$ trivial
  over k (only $\mathrm{MR}_{\mathrm{eq}}$ remains), so subscript k degenerates.
\item
  Denominator $|\mathrm{mut}_j(S_i)|$: under L5 $\wedge$ L6 $\to$
  $|\mathrm{mut}_j^{\mathrm{syntax}}(S_i)|$ (Lemma G.3).
\item
  Denominator $|\mathrm{equiv}_{i,k,j}|$: under L1 $\wedge$ L2 $\to$
  $|\mathrm{equiv}_{i,j}^{\mathrm{classic}}|$ (Lemma G.1;
  almost-everywhere w.r.t. $D_S$).
\end{itemize}

Substituting into the SMS formula gives $\mathrm{MS}_{i,j}$.

\textbf{Corollary G (LRCA trivialisation).} Under $L = L1 \wedge L2 \wedge L3$, C
= \{C1, \ldots, C5\} degenerates to \{C1\}.

\emph{Sketch.} Each of C2-C5's triggering precondition depends on at
least one $L_j$ being violated. When L1 $\wedge$ L2 $\wedge$ L3 hold simultaneously,
every non-trivial-space dimension (MR non-triviality, AVP tolerance
non-zero, non-empty equiv set, MR-design DOF, class-mapping openness)
closes; C2-C5 triggering set becomes empty (read off from §A.2 decision
tree). The per-C\_k to per-L\_j minimum-sufficient mapping depends on
the main manuscript's LRCA-diagnosis engineering thresholds; we do not claim one-to-one
correspondence at formal level, readers can trace through §A.2 + §C.4.
Under L, $\texttt{suspect\_share} \to 0$, LRCA reports only C1, SMS degenerates to
single-layer metric consistent with \citet{jia2011analysis} MS engineering
structure.

\textbf{Empirical consistency.} Theorem G and the LRCA-trivialisation
corollary jointly
show that the SMS-based empirical conclusions (Cliff's delta in the aligned-versus-cross results,
Friedman chi\^{}2 in the cross-class consistency results, and Spearman rho in the SMS-vs-Pattern-Coverage analysis) are structurally consistent
with existing \citet{jia2011analysis} literature in classical
syntactic-mutation scenarios, no metric-level semantic fragmentation.

\section{Adjoint Extension Study}\label{h.-adjoint-extension-arm}

This appendix gives the full specification of the confirmatory adjoint
extension arm summarised in the main manuscript. The arm
exercises the adjoint invariant $\psi_6$, the MetaPattern
$m_{\mathrm{adj}}$ (self-adjoint block $T^{*}$) in NOETHER's taxonomy
\citep{noether2026}, on three third-party
linear-operator libraries, disjoint from the 12 PUTs and the
pre-registered 60-cell. Each library carries a real fixed defect in its
adjoint code path as a ground-truth anchor:

\begin{itemize}

\item
  \textbf{scipy.sparse.linalg LinearOperator} (linear algebra; BSD-3).
  Program under test: the scaled-operator adjoint $M^{\mathrm{H}}$ for
  $M = \alpha A$ with complex scale $\alpha$. Real anchor: scipy issue
  \#8900 / PR \#8962 (the adjoint dropped the conjugate on $\alpha$,
  returning $\alpha A^{\mathrm{H}}$ instead of $\bar{\alpha}
  A^{\mathrm{H}}$).
\item
  \textbf{pylops.signalprocessing.FFT} with \texttt{real=True} (signal
  processing; BSD-3). Program under test: the real-FFT adjoint
  \texttt{\_rmatvec}; the adjoint identity is the library's own
  \emph{dot-test}. Real anchor: pylops issue \#268 / PR \#292 (commit
  \texttt{bd3a919}), a wrong scaling of the middle/Nyquist frequency bin
  that made the adjoint inconsistent with the forward.
\item
  \textbf{jax} \texttt{linear\_transpose} / \texttt{vjp} (automatic
  differentiation; Apache-2.0). Program under test: the autodiff-derived
  transpose $M^{\mathrm{T}}$ of a linear operator containing an inverse
  real FFT, $M(x) = \mathrm{irfft}(\mathrm{rfft}(x)\cdot w)$. Real anchor:
  jax issue \#6223 / PR \#6232 (commit \texttt{112ae41}), an
  \texttt{irfft} transpose rule that reversed order with an off-by-one
  index, so \texttt{vjp} silently produced a non-adjoint transpose.
\end{itemize}

For each subject the MR families are adj.self (self-adjointness) and
adj.dual (scaled/dual-operator transpose). The mutant pool spans
adjoint-breaking fault mechanisms (drop-conjugate, sign flip, scale and
phase errors, transpose-without-conjugation, Nyquist mis-scaling), with
the real fixed defect included among them, plus semantically equivalent
mutants (scalar-multiplication reorder, conjugation of a real adjoint
output) that a strong MR must \emph{not} kill. Because the pre-fix
releases are uninstallable on the host platform (no \texttt{arm64} wheel),
each real defect is reproduced in extracted-diff form on the current
release, mirroring the boundary study; $\varepsilon_{\mathrm{tol}}$ is an
explicit factor throughout.

{%
\begin{longtable}[]{@{}p{2.0cm}p{3.2cm}p{2.5cm}p{1.8cm}p{1.8cm}@{}}
\caption{Adjoint extension study, forward kill.}\label{tab:suppl-adjoint}\\
\toprule\noalign{}
Substrate & Operator class (real anchor) & Active killed / total & Equiv.\ survived & Forward SMS \\
\midrule\noalign{}
\endfirsthead
\endhead
\bottomrule\noalign{}
\endlastfoot
scipy LinearOperator & linear algebra (\#8900) & 5 / 5 & 2 / 2 & \textbf{1.00} \\
pylops FFT real & signal processing (\#292) & 4 / 4 & 2 / 2 & \textbf{1.00} \\
jax vjp / transpose & automatic differentiation (\#6223) & 4 / 4 & 2 / 2 & \textbf{1.00} \\
\end{longtable}
}

Across all three substrates the correct program survives, every
non-equivalent (active) mutant is killed, including the three real
fixed defects, and no semantically equivalent mutant is killed, so the
forward SMS is 1.00 with zero false positives on the equivalent set. The
scipy \#8900 defect is shared with the strong-boundary study in the main manuscript (adj.self case)
and is not counted as independent corroboration twice; the pylops \#292
and jax \#6223 defects are unique to this arm. The pools are
hand-constructed and low-cardinality, so the arm demonstrates the duality
forward rather than measuring a high-adequacy MR set on a discovered fault
population (main-paper adjoint-arm reproduction-fidelity threat).

\section{Real-Defect Evidence Summary}\label{i.-result-level-real-defect-evidence-ledger}

This appendix records only the result-level facts used by the main
manuscript's real-defect evidence slice. It deliberately excludes
candidate mining, rejected cases, repository tooling, and benchmark-release
narrative. The purpose is auditability for the semantic-mutation argument:
kill-rate, semantic-stratum alignment, and real-defect detection are related
but distinct constructs.

{%
\small
\begin{longtable}[]{@{}p{0.18\linewidth}p{0.52\linewidth}p{0.24\linewidth}@{}}
\caption{Result-level real-defect evidence summary.}\label{tab:realdefect-ledger}\\
\toprule\noalign{}
Audit item & Evidence content used in the manuscript & Main-text role \\
\midrule\noalign{}
\endfirsthead
\endhead
\bottomrule\noalign{}
\endlastfoot
Verified stratum counts & 34 verified MR-detectable defects across five
NOETHER MetaPatterns \citep{noether2026}: \(m_{\mathrm{mono}}\) 10,
\(m_{\mathrm{conv}}\) 9,
\(m_{\mathrm{inv}}\) 6, \(m_{\mathrm{adj}}\) 5, and
\(m_{\mathrm{rev}}\) 4. This real-defect corpus spans a \emph{different}
block subset than the MP1--MP5 experimental grid: it adds
\(m_{\mathrm{adj}}\) and \(m_{\mathrm{rev}}\) and omits \(m_{\mathrm{dyn}}\)
(MP4) and \(m_{\mathrm{cmp}}\) (MP5), so these five are not the MP1--MP5
strata. & Shows that the declared semantic strata are not
only toy labels \\
Mutation-phase aggregate & The semantic-mutation track exceeds the
closest baseline by +0.101 on mean score (Holm-adjusted \(p = 0.046\),
Cliff's \(\delta = +0.247\)); the remaining pairwise contrasts are not
significant & Supports a modest, qualified
mutation-score advantage rather than a dominance claim \\
Real-defect face & The semantic-mutation track detects all 34 tabled
real defects; the two generic baseline groups have nonzero detection in
7/34 and 8/34 cases, and the two ablated variants miss 19/34 and 17/34
cases, respectively, with 11 shared misses & Separates real-defect detection from aggregate mutant
kill-rate \\
Non-nesting cases & Four verified non-nesting defects (A-LAPACK-004,
A-OPENBLAS-001, B-POCKETFFT-002, and E-ORDINARYDIFFEQ-001) spanning
linear
algebra, FFT, and ordinary-differential-equation solvers & Supports the semantic-effect fiber view rather
than a single scalar hierarchy \\
\end{longtable}
\normalsize
}

This summary is intentionally a result table, not a corpus or benchmark
description. The main manuscript uses it only to support construct
separation: a relation family can score similarly on aggregate mutant
kill-rate while differing sharply on real-defect detection, and different
relation families can observe non-nested semantic-effect fibers.

\section{Study 2 Confirmatory Audit: Per-Family, Per-Class, and Incident Detail}\label{j.-study2-confirmatory-detail}

This appendix records the per-family and per-class detail behind the
Study-2 confirmatory scoreboard in the main text, the calibration and
confirmatory incident log, and the review-packet protocol. The
registration of record is
\texttt{PREREGISTRATION\_STUDY2\_v1.1.md} and the append-only pilot log is
\texttt{PILOT\_LOG.md}; both are archived with the reproducibility
package. Every number below traces to a pre-frozen SSOT.

\subsection{Operator instantiability (H1\('\))}\label{j.1-h1-instantiability}

For each operator family, Table~\ref{tab:study2-h1} gives the count of
the 28 confirmatory PUTs on which the family produced \(\geq 5\)
non-equivalent admitted mutants, against the outcome-independent coverage
ceiling from the operator registry. The registered bar is \(\geq 4\) of 5
families clearing \(\geq 8\) of 28 PUTs; all five clear. Source SSOT
\texttt{h1\_instantiability\_v5.json}. Pilot PUTs \(\{a2, b4\}\) are
excluded.

{%
\small
\begin{longtable}[]{@{}p{0.16\linewidth}p{0.20\linewidth}p{0.16\linewidth}p{0.16\linewidth}p{0.20\linewidth}@{}}
\caption{Study-2 H1\('\) per-family instantiability on the 28
confirmatory PUTs.}\label{tab:study2-h1}\\
\toprule\noalign{}
Family & Coverage ceiling & PUTs cleared & Clears $\geq 8$? & Note \\
\midrule\noalign{}
\endfirsthead
\endhead
\bottomrule\noalign{}
\endlastfoot
CE & 23 & 23 & yes & clears every applicable PUT \\
OS & 14 & 14 & yes & clears every applicable PUT \\
HP & 21 & 21 & yes & clears every applicable PUT \\
TF & 15 & 15 & yes & clears every applicable PUT \\
SI & 10 & 8 & yes & narrow family; registration expected a possible miss, clears at 8 \\
\midrule
Total & --- & 5/5 & --- & bar is 4/5; verdict CONFIRM \\
\end{longtable}
\normalsize
}

\subsection{Cross-class direction consistency (H3\('\))}\label{j.2-h3-class}

Table~\ref{tab:study2-h3} gives the class-mean aligned and cross SMS and
the resulting direction for each of the four design classes, with the
descriptive (under-powered, non-confirmatory) per-class binomial
sign-test \(p\). Three of four classes are positive; class C (surrogate
models) is negative and reported as such. Source SSOT
\texttt{h3\_class\_consistency\_v5.json}. The across-meta-pattern Friedman
statistic is \(\chi^2 = 26.23\) over 28 blocks and 5 treatments,
exploratory only.

{%
\small
\begin{longtable}[]{@{}p{0.10\linewidth}p{0.08\linewidth}p{0.20\linewidth}p{0.20\linewidth}p{0.14\linewidth}p{0.14\linewidth}@{}}
\caption{Study-2 H3\('\) per-class direction.}\label{tab:study2-h3}\\
\toprule\noalign{}
Class & $n$ PUTs & Mean aligned SMS & Mean cross SMS & Direction & Sign-test $p$ (desc.) \\
\midrule\noalign{}
\endfirsthead
\endhead
\bottomrule\noalign{}
\endlastfoot
A & 7 & 0.3333 & 0.0119 & positive & 0.0312 \\
B & 6 & 0.2222 & 0.0903 & positive & 0.1875 \\
C & 7 & 0.0476 & 0.0714 & negative & 0.9688 \\
D & 8 & 0.3333 & 0.0937 & positive & 0.0352 \\
\midrule
\multicolumn{6}{@{}p{0.95\linewidth}@{}}{3/4 classes positive; bar is $\geq 3/4$; verdict CONFIRM. Per-class $p$ values are descriptive only.} \\
\end{longtable}
\normalsize
}

\subsection{Attribution purity (H4\('\))}\label{j.3-h4-attribution}

Table~\ref{tab:study2-h4} gives the per-family multi-stratum leakage
underlying the H4\('\) miss. The mean \texttt{suspect\_share} over the
140 confirmatory cells is 0.1714, above the registered 0.05 bar. CE and
HP stay clean; the leakage concentrates in TF and OS, with CF and SI
contributing the remainder. Source SSOT \texttt{s5\_purity\_v5.json};
LRCA machinery is imported unchanged from the Study-1 S5 audit.

{%
\small
\begin{longtable}[]{@{}p{0.14\linewidth}p{0.22\linewidth}p{0.24\linewidth}p{0.30\linewidth}@{}}
\caption{Study-2 H4\('\) per-family multi-stratum leakage.}\label{tab:study2-h4}\\
\toprule\noalign{}
Family & Mutants scored & Multi-stratum & Note \\
\midrule\noalign{}
\endfirsthead
\endhead
\bottomrule\noalign{}
\endlastfoot
CE & 207 & 0 & clean \\
HP & 189 & 0 & clean \\
OS & 123 & 27 & newly leaks (0 in Study 1) \\
TF & 135 & 72 & leaks most; single-stratum spec not contained on widened grid \\
CF & 18 & 9 & residual leakage \\
SI & 69 & 9 & residual leakage \\
\midrule
\multicolumn{4}{@{}p{0.95\linewidth}@{}}{Mean \texttt{suspect\_share} 0.1714 $>$ 0.05 bar; verdict NOT\_CONFIRMED; offending families CF/OS/SI/TF.} \\
\end{longtable}
\normalsize
}

\subsection{Incident and deviation log}\label{j.4-incident-log}

Table~\ref{tab:study2-incidents} summarises the calibration-pilot and
confirmatory incidents. The pilot phase (P1--P5) fixed code defects only,
under the registration firewall (no threshold, estimand, DGP,
primary-meta-pattern, or roster change). The confirmatory phase carried
two tooling incidents (P6, P7) with no protocol impact, one disclosed
script-interface deviation (D-A1), and one disclosed design-intent defect
(P8) that leaves the registered decision rule intact but was not realised as
intended.

{%
\small
\begin{longtable}[]{@{}p{0.08\linewidth}p{0.20\linewidth}p{0.66\linewidth}@{}}
\caption{Study-2 incident and deviation log.}\label{tab:study2-incidents}\\
\toprule\noalign{}
\# & Class & Summary \\
\midrule\noalign{}
\endfirsthead
\endhead
\bottomrule\noalign{}
\endlastfoot
P1 & Near-miss (recovered) & Pool-build tool ignored CLI arguments and wiped the frozen Study-1 v3 pools from an empty cache; recovered by \texttt{git restore}; no artefact lost, no commit reached. \\
P2--P3 & Code defect (pilot) & Pool builder gained argument parsing, a v5 version map (N=30, cross-source cache), and four delete-safety guards (empty-cache refusal, no empty-overwrite, wrong-version and frozen-version deletion guards). \\
P4 & Code defect (pilot) & Response filename convention not stated in packets; an explicit \texttt{response\_filename} field added to generation/review/arbitration packets. \\
P5 & Code defect (pilot) & E1$\wedge$E2 equivalence semantics not defined in-packet; operational definitions embedded verbatim in the blinded review and arbitration instructions. \\
P6 & Tooling incident & Scoring tool defaulted to the Study-1 track and overwrote \texttt{sms\_track1.json}; restored before use and rerun with an explicit track-2 target. Deterministic re-measurement of frozen pools, not regeneration. \\
P7 & Blinding guard fired & One PUT's mutants carried a source-family token in docstrings; presentation-layer redaction added at export (displayed code only), affected packets re-exported, assertion retained as the final guard. \\
D-A1 & Disclosed deviation & Pre-frozen $\delta$ script gained a post-data mode that skips only the gated H2-2 computation and emits its registered not-run verdict; H2-1\('\) logic byte-unchanged; recorded in the output SSOT. \\
P8 & Disclosed design-intent defect & The CF/TF single-stratum admission screen was wired but a silent no-op for the whole v5 campaign: the pool builder tagged op-ids with a model-source suffix (e.g.\ \texttt{c7\_TF1\_claude}) that the screen's category regex rejects, so category lookup returned \texttt{None} and every candidate was admitted without a flip-count check (81/81 CF/TF null-pass; the pilot's ``9/9 pass'' was the same false negative). H4\('\) thus evaluated the \emph{unscreened} pool. The registered decision rule scores \texttt{suspect\_share} on the admitted pool and does not condition on the screen, so the frozen NOT\_CONFIRMED verdict is valid; only the design intent (CF/TF entering single-stratum) was not realised. The filter code is deliberately \emph{not} fixed here; correcting it belongs to the Study-3 registration wave. See the post-hoc diagnosis in the main text. \\
\end{longtable}
\normalsize
}

\subsection{Review-packet protocol}\label{j.5-packet-protocol}

Confirmatory review used the same dual-blind pipeline as the registered
protocol: generation, then blind review on packets that expose only the
mutant code, the operator spec, and the PUT source (generator identity,
arm label, and SMS withheld), then arbitration on disagreement, then a
freeze-then-score step in which SMS is computed only after review labels
are frozen. The harness spawns each reviewer and the arbiter as separate,
role-isolated Claude instances, so generator, reviewer, and arbiter
differ on every item; all instances share one model family (same-vendor).
The confirmatory pools' per-mutant source tokens
\texttt{claude}/\texttt{gpt}/\texttt{deepseek} are nominal slot labels
inherited from the Study-1 slot layout, all served by the same
Claude-family generator, not real GPT or DeepSeek outputs; the
aligned-versus-cross split is meta-pattern-based and independent of these
slot labels, so no confirmatory statistic partitions by the vendor-named
slot.
In the confirmatory run, 18 independent reviewer instances returned 747
verdicts, 4 mutants were rejected at the first round, and 6 uncertain
cases were arbitrated by a third instance and all resolved as confirmed
under one uniform rationale (pinning the Gaussian-process hyperparameters
before comparison). Mechanical admission gates V1--V4 admitted 756 of 774
valid mutants, yielding 140 confirmatory cells (36 with nonzero killed
mass). Because the harness is same-vendor, the cross-vendor dual-blind
hypothesis H2-2 is gated not-run, with no same-vendor proxy substituted.

\subsection{Study-3 graded-attribution follow-up}\label{j.6-study3-detail}

Study 3 re-registers the attribution construct alone under
\texttt{PREREGISTRATION\_STUDY3\_v2.md} (v2.0); the append-only pilot log is
extended in the Study-3 section of \texttt{PILOT\_LOG.md}. It replaces the
single Study-2 purity bar with two Holm(2) verdicts in Family G and re-opens no
other Study-2 verdict. The confirmatory statistics trace to the pre-frozen SSOTs
\texttt{sms\_track2\_v6.json} (fresh validated pool) and
\texttt{h4\_graded\_v6.json} (both verdicts); the power reference is
\texttt{power\_study3.json}. Pilot PUTs \(\{a2, b4\}\) are excluded from every
confirmatory statistic.

Table~\ref{tab:study3-exec} records the one-shot execution, and
Table~\ref{tab:study3-verdicts} the two registered verdicts with their per-class
and per-family breakdowns.

{%
\small
\begin{longtable}[]{@{}p{0.46\linewidth}p{0.48\linewidth}@{}}
\caption{Study-3 confirmatory execution (one shot, seed 20260708).}\label{tab:study3-exec}\\
\toprule\noalign{}
Stage & Value \\
\midrule\noalign{}
\endfirsthead
\endhead
\bottomrule\noalign{}
\endlastfoot
Valid generation records & 765 \\
Admitted (gates V1--V4, live all-family screen) & 720 \\
Rejected at admission & 45 (families a3, b3, c1, d8) \\
v6 confirmatory pool & 633 mutants over 28 PUTs \\
v6 pool including \(\{a2, b4\}\) pilot & 678 files over 30 directories \\
Confirmatory cells & 140 (20 with nonzero killed mass) \\
Blind-review verdicts & 720 (0 schema errors, 0 arbitrations) \\
Live-screen re-audit & 633 candidates matched, 36 flagged multi-stratum
(v5 no-op matched 0) \\
\end{longtable}
\normalsize
}

{%
\small
\begin{longtable}[]{@{}p{0.13\linewidth}p{0.20\linewidth}p{0.14\linewidth}p{0.45\linewidth}@{}}
\caption{Study-3 Family-G verdicts (Holm 2).}\label{tab:study3-verdicts}\\
\toprule\noalign{}
Hypothesis & Threshold & Verdict & Breakdown \\
\midrule\noalign{}
\endfirsthead
\endhead
\bottomrule\noalign{}
\endlastfoot
H4\('' \)-strict & purity $\geq 0.90$ (95\% lower CP) & CONFIRM & purity 1.0,
CP-lower 0.9673 over $n = 90$ detected clean-family mutants (CE 72, HP 18, CF 0);
screen matched $>0$, gate passes \\
H4\('' \)-graded & rich share $\geq 0.15$ (95\% lower boot) & NOT\_CONFIRMED &
rich-class mean share 0.0833, lower bound 0.0, $n_{\mathrm{rich}} = 6$
(C3--C6, D2, D5); per-class C 0.0 (4 PUTs), D 0.25 (2 PUTs) \\
\end{longtable}
\normalsize
}

\textbf{Live screen, P8 remediation verified.} The P8 fix (op-id parser
tolerant of the model-source suffix, all-family scope, loud-fail on a null
category) is exercised in the confirmatory campaign, not only in regression
tests. The all-family re-audit over the 633-mutant pool matched all 633
candidates and flagged 36 as multi-stratum
(\texttt{n\_screened\_candidates} 633, \texttt{n\_multistratum\_flagged} 36 in
\texttt{h4\_graded\_v6.json}); the byte-identical v5 screen matched zero, so the
registered screen-smoke gate, which requires a positive match, passes loudly
rather than admitting an unconstrained pass.

\textbf{Defect P9 (only post-freeze code change).} The pilot surfaced one
code-level defect: \texttt{build\_pools} and \texttt{sms\_campaign} knew only
pool versions v2--v5. The fix, mirroring the v5 D2/D4 wiring, adds a \texttt{v6}
version (suffix \texttt{\_pool\_v6}, cross-source cache, N = 30, strict pool
resolution) and builds v6 under the live all-family screen; a wrong-version
deletion guard means a v6 build can only ever remove a \texttt{\_pool\_v6}
directory, so the frozen v5 pools are never touched. It encodes the registered N
for a new pool version and points scoring at the correct strict directory; it
changes no threshold, estimand, primary-meta-pattern assignment, or roster, and
reads no Study-3 outcome.

\bibliographystyle{ACM-Reference-Format}
\bibliography{references}

\end{document}
